\documentclass{article}
\usepackage[utf8]{inputenc}
\usepackage[T1]{fontenc}
\usepackage{geometry}
\usepackage{graphicx}
\usepackage{float}
\usepackage{hyperref}
\usepackage[english,indonesian]{babel}

\title{Optimasi Pemrosesan Template LaTeX Menggunakan Pendekatan Berbasis AI}
\author{Academic Support Agent}
\date{\today}

\begin{document}
\maketitle

\begin{abstract}

\end{abstract}

\chapter{PENDAHULUAN}


\chapter{TINJAUAN PUSTAKA}


% Original content (commented out):
% %Nicht anfassen, so ist der Dokumentenaufbau
\documentclass[a4paper, 12pt]{article}
\input{settings/packages}
\input{settings/page}
\begin{document}
\input{settings/titlepage}
\setcounter{page}{1}
%%%%%%%%%%%%%%%%%%%%%%%%%%%%%
%%%Ab hier Inhalt einfügen%%%
%%%%%%%%%%%%%%%%%%%%%%%%%%%%%

\section{Einleitung} \label{einleitung}

Text der Einleitung


\section{Hauptteil}

Text des Hauptteils

\subsection{Unterüberschrift 1}

Hier wird der Artikel von \textcite{Abadi.2016} zitiert

\section{Schluss}


%%%%%%%%%%%%%%%%%%%%%%%%%%%%
%%%%%Ende des Dokuments%%%%%
%%%%%%%%%%%%%%%%%%%%%%%%%%%%
\newpage
\input{settings/endofdocument}
\end{document}

\end{document}